\documentclass{vii-coloquio-co}
\usepackage[utf8]{inputenc}
\usepackage[T1]{fontenc}
\usepackage{geometry}
\usepackage{graphicx}
\usepackage{float}
\usepackage{hyperref}
\usepackage[english,indonesian]{babel}

\title{Optimasi Pemrosesan Template LaTeX Menggunakan Pendekatan Berbasis AI}
\author{Academic Support Agent}
\date{\today}

\begin{document}
\maketitle

\begin{abstract}

\end{abstract}

\chapter{PENDAHULUAN}


\chapter{TINJAUAN PUSTAKA}


% Original content (commented out):
% % ----------------------------------------------------------------------------------- %
% VII Col\'oquio de Matem\'atica da Regi\~ao Centro-Oeste
% UFGD - Dourados, MS - 08 a 12 de abril de 2024
% 
% O modelo fornece dois comandos para auxiliar ne declara\c{c}\~ao de figuras.
% 
% Uma \'unica figura:
% \figura{fator_de_escala}{nome_do_arquivo}{legenda_da_figura}
%
% Duas figuras lado a lado ocupando toda a largura da linha:
% \figuras{fator_de_escala}{arquivo_da_figura1}{arquivo_da_figura2}{legenda_da_figura}
%
% O comando \colunas organiza o conte\'udo do p\^oster em duas colunas.
% ----------------------------------------------------------------------------------- %

\documentclass{vii-coloquio-co}

\usepackage{lipsum}  % utilizado para gerar o texto do modelo

% declara\c{c}\~ao do t\'{\i}tulo e autores do trabalho
\titulo{T\'{\i}tulo bem criativo}
\autor{Primeiro Autor*, Segundo Autor**}
\email{*primeiro@autor.edu.br, **segundo@autor.edu.br}
\universidade{Universidade Federal da Grande Dourados}

\begin{document}

\cabecalho

\section*{Resumo}
\lipsum[1]

\section{Primeira se\c{c}\~ao}
\lipsum[1-2]

\figura{1.0}{figura.png}{Figura com tamanho original.}

\lipsum[3-4]\cite{texbook}

\section{Segunda se\c{c}\~ao}
\lipsum[1-2]\cite{latex2e}

\figura{2.0}{figura.png}{Figura com o dobro do tamanho.}
\figura{0.5}{figura.png}{Figura com metade do tamanho.}
\figuras{figura.png}{figura.png}{Figuras lado a lado.}

\lipsum[3-4]

\section{Conclus\~ao}
\lipsum[1]

\section*{Agradecimentos}
\lipsum[4]

\bibliographystyle{plain}
\bibliography{referencias}

\end{document}


\end{document}